% reproducibility.tex
\label{sec:reproducibility}

\begin{itemize}
\item \textbf{Environment.}  Python~3 with \texttt{numpy},
      \texttt{scipy}, \texttt{matplotlib}, \texttt{pytest}.
      Exact dependency list in
      \texttt{virtual-env-requirements.txt}.  The \texttt{setup.sh}
      / \texttt{setup.cmd} scripts create a \texttt{.venv} and
      install everything.

\item \textbf{Engine.}  The experiment depends on the A=1 framework
      engine via the shared \texttt{dcl\_core} package
      (\texttt{dcl\_core.core.OctahedralLattice},
      \texttt{dcl\_core.core.CausalSession}, \dots).
      \texttt{dcl\_core} is declared in
      \texttt{virtual-env-requirements.txt} and installed by
      \texttt{setup.cmd} / \texttt{setup.sh} automatically; no
      separate engine setup step is required.  See \texttt{CLAUDE.md}'s
      \emph{Engine} section for the architectural rationale and the
      pinning strategy at release time.

\item \textbf{Build the paper.}  From a fresh clone, with the venv
      activated:
      \begin{verbatim}
./build.sh paper        # POSIX / MSYS2 UCRT64 on Windows
build.cmd paper         # Windows cmd / PowerShell
      \end{verbatim}
      runs \texttt{pdflatex} three times and \texttt{bibtex} once,
      producing \texttt{paper/main.pdf}.

\item \textbf{Run the experiment.}  Once the engine is provisioned:
      \begin{verbatim}
python -u src/experiments/exp_18_tidal_ionization.py \
    > data/exp_18_tidal_ionization.log 2>&1
      \end{verbatim}
      Wall-clock runtime: a few hours on a modern laptop (7
      gradient values $\times$ up to 2500 ticks each, with each
      tick a 3D wavefunction step on a $65^3$ grid).  Expected
      output: monotonic $T_\text{escape}(g)$ decrease and an
      opposite-sign tidal-displacement signature
      ($x_e\_proj \cdot x_p\_proj < 0$) for $g > 0$.

\item \textbf{Master audit roll-up.}  From the repo root:
      \begin{verbatim}
python audit_universe.py
      \end{verbatim}
      reports the declared status of every row in
      \texttt{paper/sections/audit\_table.tex} and parses the
      latest \texttt{data/exp\_18*.log} for the actual cached
      \texttt{PASS}/\texttt{FAIL} marker.

\item \textbf{Where to look next.}  Start with
      \texttt{notes/quantum\_roche\_limit.md} for the analytical
      derivation; then
      \texttt{notes/plasma\_as\_gravitational\_ionization.md}
      (inherited from Paper~I) for the broader gradient-ionization
      theory; then the experiment's \texttt{.md} companion for
      runtime parameters and audit criteria.
\end{itemize}
